% The rate without the extent (paper 11 of the O lane).
% Lane: O (operator bridge), module YangMills/OS/SpatialUniform.lean.
% REGIME: tricotomy; NO scores, NO desk language, NO commit chronology beyond
% the reproducibility freeze.
\documentclass[11pt]{article}
\usepackage[margin=1.1in]{geometry}
\usepackage{amsmath,amssymb,amsthm}
\usepackage{booktabs,longtable,array}
\usepackage[colorlinks=true,linkcolor=blue,urlcolor=blue,citecolor=blue]{hyperref}

\newtheorem{theorem}{Theorem}[section]
\newtheorem{lemma}[theorem]{Lemma}
\newtheorem{proposition}[theorem]{Proposition}
\newtheorem{corollary}[theorem]{Corollary}
\newtheorem{remark}[theorem]{Remark}
\theoremstyle{definition}
\newtheorem{definition}[theorem]{Definition}

\emergencystretch=3em
\tolerance=2000

\newcommand{\anchor}{af88cd6abc95c871495f7ac911690e938dc177b3}
\newcommand{\osline}[3]{%
  \href{https://github.com/lluiseriksson/THE-ERIKSSON-PROGRAMME/blob/\anchor/YangMills/OS/#1.lean\#L#2}{\nolinkurl{#3}}}
\newcommand{\repofile}[2]{%
  \href{https://github.com/lluiseriksson/THE-ERIKSSON-PROGRAMME/blob/\anchor/#1}{\nolinkurl{#2}}}
\newcommand{\R}{\mathbb{R}}
\newcommand{\Z}{\mathbb{Z}}
\newcommand{\SU}{\mathrm{SU}}

\title{The Rate Without the Extent:\\
a Machine-Checked Uniform Spectral Modulus\\
for the Decoupled $\Z_2$ Slice\\
\large and a pre-registered judge that failed at the coupled one}
\author{Lluis Eriksson}
\date{}

\begin{document}
\maketitle

\begin{abstract}
Every rate in this lane so far has been a fixed-extent rate.  The gap paper
proved strict spectral separation at each extent and said plainly that it was
not uniform \cite{ErikssonGap}; the modulus paper gave that separation a
number, $\mathrm{specRatio}(L)$, and reported measurements saying the number
tends to $1$ outside the disordered region \cite{ErikssonSpectral}.  A geometric
bound whose rate tends to $1$ is empty in the volume limit, so nothing in the
lane survived $L\to\infty$, and the word \emph{clustering} was never used.

\textbf{Proved.}  For the \emph{decoupled} kernel --- the transfer kernel at
constant source weight --- the modulus is
$\mathrm{specRatio} = \tanh\beta$ at \emph{every} extent, with $L$ nowhere in
it.  Both directions: an operator bound by induction on the extent, and
attainment by the single-site observable the extent paper already built
\cite{ErikssonExtent}.  Composing with the modulus paper's endpoint, the
normalised Gibbs two-point function obeys
$|\mathbb{E}[A(X_0)A(X_N)]| \le C_A\,(\tanh\beta)^N$ past one threshold serving
every observable --- a bound whose \emph{rate} contains no $L$, and therefore
the first statement in this lane that survives the volume limit.

\textbf{Why the proof is not the spectral decomposition.}  The decoupled kernel
is a product over sites, so its spectrum is a product; that route needs the
spectrum of a Kronecker power, which the library does not carry.  It is not
needed.  The modulus paper proved that $\mathrm{specGap}$ is the \emph{greatest}
norm ratio on the fluctuation sector, so bounding it above is an operator
inequality and nothing else, and the operator inequality falls to induction: the
even part of an observable keeps its mean zero and inherits the rate, the odd
part keeps nothing and gets only Schur's test, and the two recombine
\emph{exactly}, because $\tanh\beta \cdot Z = D$ with $Z$ the row sum and $D$
the odd eigenvalue of a single bond.

\textbf{Not proved, and a judge that failed.}  The \emph{coupled} kernel is
untouched.  Before any of this was written we pre-registered two falsifiable
predictions.  The first --- that the decoupled rate is exactly $\tanh\beta$ at
every extent --- passed to $10^{-16}$, and authorised the work above.  The
second --- that the coupled uniformity boundary is the Onsager curve --- failed
on one of eight pre-registered cells, and it stays failed: that claim is
therefore reported as \textbf{not established}, not softened.  At constant
source weight the spatial slices are independent, so what is proved here is a
statement about a product measure; that is exactly why it is reachable, and it
is said here rather than left to be noticed.  Reflection positivity is
untouched, and nothing in this paper is a claim about $\SU(N)$, the continuum
limit, or the Yang--Mills mass gap \cite{Clay}.
\end{abstract}

\section{Ten papers of fixed-extent rates}\label{sec:debt}

The lane's rates have all been local in the extent.  That is not a stylistic
complaint: a rate $r(L)$ with $r(L)\to 1$ gives, at fixed $N$, a bound that
becomes vacuous as the system grows, and the infinite-volume statement one
actually wants --- exponential decay of correlations with a fixed rate --- is
not implied by any amount of fixed-$L$ information.

So the honest question is whether \emph{some} kernel in this family has a rate
that does not move with $L$.  This paper answers it for the decoupled kernel,
and the answer is exact.

\section{The construction}\label{sec:constr}

Let $Z = e^{\beta}+e^{-\beta}$ and $D = e^{\beta}-e^{-\beta}$, and let $K$ be the
decoupled kernel on $\{0,1\}^L$: a product over sites of the single-bond matrix.
Its row sums are all $Z^L$, so the constant observable is its Perron vector with
eigenvalue $\lambda = Z^L$.

\begin{lemma}[Schur's test]\label{lem:schur}
A symmetric nonnegative kernel all of whose row sums equal $R$ satisfies
$\|Ku\| \le R\,\|u\|$.
\end{lemma}

\osline{SpatialUniform}{83}{norm\_act\_le\_of\_rowSum}.
One Cauchy--Schwarz per row, then summing.  This is the only bound the odd part
of the induction is given, and it is all it needs.

\begin{theorem}[the uniform operator bound]\label{thm:uniform}
For every extent $L$ and every observable $u$ with $\sum_\sigma u(\sigma)=0$,
\[
  \|Ku\| \;\le\; \tanh\beta \cdot Z^{L} \cdot \|u\| .
\]
\end{theorem}

\osline{SpatialUniform}{318}{spatialKernel\_fluct\_bound}, from the squared
form \osline{SpatialUniform}{176}{spatialKernel\_fluct\_sq}, which is what the
induction runs on --- a square root in the inductive hypothesis would be a
square root at every step.

\section{Why the constant does not degrade}\label{sec:exact}

Split an observable at the first site, $u = (u_0,u_1)$, and set $p = u_0+u_1$,
$m = u_0-u_1$.  Then $p$ still has mean zero, so the inductive hypothesis
applies to it with rate $\tanh\beta$; $m$ has no orthogonality left and gets
only Lemma~\ref{lem:schur}, with rate $1$.  Acting with the kernel and
collecting,
\[
  \|Ku\|^{2} \;=\; \tfrac12\bigl(Z^{2}\|Kp\|^{2} + D^{2}\|Km\|^{2}\bigr),
  \qquad
  \|u\|^{2} \;=\; \tfrac12\bigl(\|p\|^{2}+\|m\|^{2}\bigr).
\]
Feeding in the two bounds gives
$\tfrac12 Z^{2}(\tanh\beta\,Z^{L})^{2}\|p\|^{2}
 + \tfrac12 D^{2}(Z^{L})^{2}\|m\|^{2}$, and the two coefficients are
\emph{equal}, because
\[
  \tanh\beta \cdot Z \;=\; D .
\]

\osline{SpatialUniform}{167}{tanh\_mul\_z2Norm}.
That single identity is the whole reason the constant survives: the rate given
up on the odd part is exactly the odd eigenvalue of one bond, at every step.
Lose it and the constant would pick up a factor per site.

\section{Sharpness}\label{sec:sharp}

Theorem~\ref{thm:uniform} is an inequality.  It is attained.

\begin{theorem}[the rate is exactly $\tanh\beta$]\label{thm:sharp}
At every extent with at least one site,
$\mathrm{specGap} = \tanh\beta \cdot Z^{L}$, hence
$\mathrm{specRatio} = \tanh\beta$.
\end{theorem}

\osline{SpatialUniform}{404}{spatialKernel\_specGap\_eq}.
The upper bound is Theorem~\ref{thm:uniform} together with the fact that
$\mathrm{specGap}$ is the \emph{greatest} norm ratio on the fluctuation sector
\cite{ErikssonSpectral} --- which is the only place that paper is used, and it
is used as a bound, never as a spectral decomposition.  The lower bound is the
single-site sign observable of the extent paper \cite{ErikssonExtent}: it has
mean zero (\osline{SpatialUniform}{376}{sum\_siteSign}), it is nonzero, and
its eigenvalue $D\,Z^{L-1}$ realises the bound exactly.

\section{The payoff}\label{sec:payoff}

\begin{theorem}[a rate that survives the volume limit]\label{thm:cluster}
At constant source weight there is an $N_0$ such that for every observable $A$
whose dressing has mean zero there is a $C_A>0$ with
\[
  \bigl|\mathbb{E}[A(X_0)A(X_N)]\bigr| \;\le\; C_A\,(\tanh\beta)^{N},
  \qquad N \ge N_0 ,
\]
and $\tanh\beta$ does not depend on the extent.
\end{theorem}

\osline{SpatialUniform}{499}{gibbs\_clustering\_uniform\_rate}, composing
Theorem~\ref{thm:sharp} with the modulus paper's endpoint
\cite{ErikssonSpectral}.  The hypothesis is not vacuous: that paper's
non-vacuity lemma supplies a nonzero observable satisfying it whenever the state
space has two points.

\textbf{What is uniform and what is not.}  The \emph{rate} is uniform in $L$.
The constant $C_A$ is not claimed to be, and depends on the observable; the
threshold $N_0$ does not depend on the observable, by the modulus paper's
uniform-threshold theorem, but is not claimed uniform in $L$ either.

\section{Scope, and a judge that failed}\label{sec:scope}

\paragraph{This is a product measure.}  At constant source weight the spatial
slices decouple, so the Gibbs measure is a product across sites and the
statement above is, in substance, the one-dimensional chain repeated.  That is
why it is reachable.  What it establishes is narrower than it may sound: not
that the coupled model clusters uniformly, but that the \emph{obstruction to
uniformity is the source weight and not the extent}.

\paragraph{The coupled kernel is exactly a 2D Ising transfer matrix.}  With the
ring source weight, $K = D_w^{1/2} T D_w^{1/2}$ where $T$ is the inter-slice
product and $D_w$ the intra-slice ring weight --- which is the symmetrised
transfer matrix of the two-dimensional Ising model, vertical coupling $\beta$
and horizontal coupling $\gamma$.  This is an exact algebraic identification,
not a numerical claim.

\paragraph{And a pre-registered judge that failed.}  Before any of the above was
written we committed two falsifiable predictions with their thresholds fixed in
advance \repofile{scripts/judge\_spatial\_uniform.py}{scripts/judge\_spatial\_uniform.py}.
The first --- the decoupled rate is exactly $\tanh\beta$ at every extent ---
passed to $10^{-16}$ and authorised this paper.  The second --- that the
\emph{coupled} uniformity boundary is the Onsager curve
$\sinh 2\beta \sinh 2\gamma = 1$ --- predicted saturation below $1$ on the five
disordered cells and a rate tending to $1$ on the three ordered ones.  Seven of
eight cells agreed.  The cell $(\beta,\gamma)=(0.1,1.0)$, disordered, gave
$\mathrm{specRatio}(12)=0.6266$ --- nowhere near $1$ --- but with drift from
$L=10$ of $3.5\times10^{-2}$, above the pre-registered $10^{-2}$, so the cell is
\emph{undecided at $L\le12$} and the judge does not pass.

We did not raise the extent until the cell passed.  The autopsy
\repofile{scripts/judge\_spatial\_uniform\_autopsy.py}{scripts/judge\_spatial\_uniform\_autopsy.py}
prints the whole sequence instead: its increments decay geometrically with ratio
$0.835$, extrapolating to about $0.707$, so this is saturation not yet reached
and not a crawl to $1$.  That is a diagnosis, not a pass.  \textbf{The Onsager
identification of the boundary is therefore reported as not established}, and no
statement in this paper rests on it.

\paragraph{No claim about the physical target.}  Reflection positivity is
untouched, and nothing here concerns $\SU(N)$, the continuum limit, or the Clay
problem \cite{Clay}.  The distance to that problem is unchanged.

\section{What is verified and not proved}\label{sec:verified}

The measured $\mathrm{specRatio}(L)$ for the coupled kernel, extended past the
five points the modulus paper printed:

\begin{center}
\begin{tabular}{llll}
\toprule
$(\beta,\gamma)$ & $L=1$ & $L=12$ & increments \\
\midrule
$(0.2,0.2)$ & $0.1974$ & $0.294449$ & geometric, ratio $0.278$ \\
$(0.3,0.4)$ & $0.2913$ & $0.646895$ & geometric, ratio $0.648$ \\
$(0.1,1.0)$ & $0.0997$ & $0.626639$ & geometric, ratio $0.835$ \\
$(0.8,1.2)$ & $0.6640$ & $1.000000$ & to $1$ \\
\bottomrule
\end{tabular}
\end{center}

These are \textbf{verified} and \textbf{not proved}; no theorem depends on them.
They are printed because they are the reason the coupled case is believed hard
rather than false, and because the last row is the one the whole lane has been
reporting as its limitation.

\section{Reproducibility}\label{sec:repro}

All artifacts are at source checkpoint \texttt{\anchor} on branch \texttt{main}.
Every hyperlink was resolved against that commit, and checked to point at an
actual declaration line, before compilation.

\begin{center}
\begin{tabular}{ll}
\toprule
Quantity & Value \\
\midrule
Full core build & 8431 jobs, success \\
Oracle commands, and answers & 2689 \\
\quad distinct declarations & 2686 \\
\quad commands listed twice & 3 \\
\quad distinct, with axiom dependencies & 2663 \\
\quad distinct, axiom-free & 23 \\
Declarations in this module & 17 \\
Distinct axioms across all outputs & \texttt{propext}, \texttt{Quot.sound}, \texttt{Classical.choice} \\
Occurrences of \texttt{sorryAx} & 0 \\
Project axioms & 0 \\
Errors & 0 \\
\bottomrule
\end{tabular}
\end{center}

The module is
\repofile{YangMills/OS/SpatialUniform.lean}{YangMills/OS/SpatialUniform.lean};
the pre-registered judge and its autopsy are
\repofile{scripts/judge\_spatial\_uniform.py}{scripts/judge\_spatial\_uniform.py}
and
\repofile{scripts/judge\_spatial\_uniform\_autopsy.py}{scripts/judge\_spatial\_uniform\_autopsy.py};
verdicts are in
\repofile{docs/VERIFICATION-LEDGER.md}{docs/VERIFICATION-LEDGER.md}.

\begin{thebibliography}{9}

\bibitem{ErikssonExtent}
L.~Eriksson,
\emph{Where the Elementary Reconstruction Stops: a Machine-Checked Obstruction
at the Coupled $\Z_2$ Slice}, viXra:2607.0075, 2026.

\bibitem{ErikssonGap}
L.~Eriksson,
\emph{Strict but Not Uniform: a Machine-Checked Spectral Gap at Every Finite
Extent}, viXra, 2026.

\bibitem{ErikssonSpectral}
L.~Eriksson,
\emph{The Modulus: a Machine-Checked Operator Bound on the Fluctuation Sector of
the Coupled $\Z_2$ Slice}, viXra, 2026.

\bibitem{Onsager}
L.~Onsager,
\emph{Crystal statistics I: a two-dimensional model with an order-disorder
transition}, Phys.\ Rev.\ \textbf{65} (1944), 117--149.

\bibitem{Seneta}
E.~Seneta,
\emph{Non-negative Matrices and Markov Chains}, 2nd ed., Springer, 1981.

\bibitem{Clay}
A.~Jaffe and E.~Witten,
\emph{Quantum Yang--Mills theory}, Clay Mathematics Institute Millennium Prize
problem description, 2000.

\end{thebibliography}

\end{document}
